% Mirror of Figure 3 caption for Overleaf paste-in.
% Matches R/revision_2026/figures/fig3_regimes.R as of 2026-08-25 --
% the CURRENT 3-panel main-text Figure 3 (A: locked recovery, B: same
% shock across the feedback-strength grid, C: history-dependence index).
%
% NOTE: an earlier caption draft (from the "symptom burden" terminology
% pass) described the OLD 2x2 off/on layout (fig3_recovery_feedback.R),
% which no longer matches this figure -- confirmed with the user
% (2026-08-25) that the 3-panel history-dependence version stays as
% main-text Figure 3, so this caption was rewritten from scratch to
% describe panels A/B/C as they actually are, using "symptom activation"
% terminology throughout instead of "symptom burden."

\caption{\textbf{Feedback strength governs recovery and history-dependence in the slow--fast model.}
(A) Under the locked moderate-feedback level (\(b = 0.50\)), an acute perturbation to the slow contextual field is followed by a transient rise in mean symptom activation (\(m_t\)); trajectories recover toward baseline whether or not symptom-to-context feedback is present.
(B) Applying the same perturbation across a range of feedback strengths shows a shift from full recovery (\(b = 0\)) to settling at a stable, elevated plateau (\(b = 1.0\)), rather than returning to the pre-shock level.
(C) A history-dependence index summarizes, across the same feedback-strength range, whether late mean symptom activation depends on the system's initial state. Trajectories initialized from a low- versus a high-activation state converge under weaker feedback but remain separated once feedback strength exceeds approximately \(b = 1.0\) (history-dependent regime), while very high feedback strengths instead approach runaway growth in the slow field rather than a genuine second equilibrium. The locked main-text feedback level (\(b = 0.50\)) falls within the convergent regime.}
